\lesson{4}{Oct 12 2021 Mon (14:38:20)}{Quantum Models}{Unit 2}

Once we have determined where an electron could be in the electron cloud, it has
already moved on, and the prediction of its location is obsolete. Therefore, we
can only describe the probable locations of electrons within an atom's electron
cloud.

The probable location of electrons within an atom can be determined using the \bf{Quantum Model}.

\begin{definition}[Quantum Model]
    A mathematical model based on quantum theory that represents the probable
    state of electrons within an atom.
\end{definition}

The probable position of every electron can be described by four quantum
numbers. The first number is called the \bf{Principal Quantum Number (n)}.
It can be any whole number from one to seven.

The principal quantum number refers to the energy levels around the nucleus
where electrons are located. These energy levels are also called energy shells.

Energy levels contain subshells that are designated with letters such as

\begin{itemize}
    \item s
    \item p
    \item d
    \item f
\end{itemize}

Due to the number of electrons in most atoms, multiple subshells exist
within each energy level.

Some subshells have multiple orbitals—regions inside an
atom where electrons are mostly likely to be found. However, each orbital (no
matter the type) can only hold two electrons.

\paragraph{The $s$ orbital} The $s$ orbital is spherically shaped. Because the electrons move so quickly, they
can be located anywhere inside the sphere.

The $s$ subshell contains $1$ orbital and is denoted by a line. The electrons
are denoted as arrows. An up arrow represents an electron that has a clockwise
spin. A down arrow represents an electron with a counterclockwise spin.

\paragraph{The $p$ orbital} The $p$ orbital is dumbbell-shaped, as shown to the left.

There are $3$ of these orbitals, one on each axis.

The $p$ subshell has $3$ orbitals and is denoted by three lines. This subshell
can hold a maximum of $6$ electrons.

\paragraph{The $d$ orbital} The $d$ orbital is nondescript in shape. The $d$ subshell has five orbitals and
is denoted by $5$ lines. This subshell can hold a maximum of $10$ electrons.

\paragraph{The $f$ orbital} The $f$ orbital is nondescript in shape.

The $f$ subshell has $7$ orbitals and is denoted by seven lines. This subshell
can hold a maximum of $14$ electrons.

A configuration describes the arrangement of objects. Like a map, it shows the
location of objects and the distance between them. Similarly, an electron
configuration represents the arrangement of electrons around the nucleus of an
atom. Each element has its own electron configuration. But unlike a map, an
electron configuration doesn't show exact locations for electrons, only their
probable placement in energy levels and subshells. Remember, within these
subshells electrons are in constant motion.

All systems in nature tend to take a position or arrangement that has the lowest
possible potential energy. The lowest potential energy arrangement for an atom's
electrons, called the ground-state electron configuration, is the arrangement
that places them closest to the nucleus.
